% Section 3.5 -- the length axis, compressed. The tactic-step proxy and its
% under-recording are in Appendix D. Framing: a reusable methods result, not a
% caveat about our data. No literal numbers: macros only.

\subsection{Length must be measured on a wrap-invariant axis}
\label{sec:length-axis}

Controlling for proof length is mandatory here. Short proofs mechanically cite
fewer and more common premises, so an uncontrolled length difference can
manufacture every effect we report. The choice of length \emph{axis} is itself
a measurement question with a checkable answer, and that is the reusable part
of this section.

Two axes are available. On characters the corpora barely overlap, at an
overlap coefficient of \charOverlapSample{}. On tactic steps they nearly
coincide, at \stepOverlapCensus{}.\footnote{The two rest on different bases:
\charOverlapSample{} is computed on the standardization sample the joint
control was cut on, \stepOverlapCensus{} over the full census of both corpora.
The step-axis figure on the sample basis is \stepOverlapSample{}, so the basis
is not what separates the axes.} The disagreement is explained by how each
corpus reaches the page. Mathlib source is hand-wrapped near a column limit and
carries docstrings; a released machine proof is generated in one pass, wrapped
differently or not at all, and here carries a large volume of natural-language
commentary. Character counts absorb all of that; tactic counts do not, because
wrapping and commentary do not change how many tactics run.

Standardizing on characters is therefore substantially standardizing on source
formatting, and that is not a neutral error. On the character axis one of our
secondary metrics changed sign and the recall term of
Section~\ref{sec:reference-line} collapsed to a null, and both reverse on the
step axis. Every length control in this paper accordingly runs on tactic steps,
and character length appears once, here, as the diagnostic that exposed the
problem. It is not an alternative specification, because a specification that
measures wrapping is not a second opinion about proof size. Counting steps
needs its own instrument and the obvious one under-records, which
Appendix~\ref{app:steps} measures and Section~\ref{sec:limitations}, L2,
carries as a limitation.

The check that settles the axis is cheap. Compute the overlap coefficient on
each candidate axis before choosing one. Where the axes disagree as sharply as
they do here, the choice is load-bearing, and making it after seeing the
comparison is a researcher degree of freedom.
